\documentclass[11pt,a4paper,fontset=fandol]{ctexart}

\usepackage[a4paper,top=2.25cm,bottom=2.45cm,left=2.25cm,right=2.25cm,headheight=18pt,headsep=0.55cm,footskip=24pt]{geometry}
\usepackage{amsmath,amssymb,amsthm}
\usepackage{fancyhdr}
\usepackage{lastpage}
\usepackage{enumitem}
\usepackage{titlesec}
\usepackage{xcolor}
\usepackage{hyperref}
\usepackage{bookmark}
\usepackage[most]{tcolorbox}
\usepackage{setspace}
\usepackage{array}

\hypersetup{
  colorlinks=true,
  linkcolor=blue!50!black,
  urlcolor=blue!50!black
}

\setstretch{1.13}
\setlength{\parindent}{2em}
\setlength{\parskip}{0.25em}
\allowdisplaybreaks
\setlist[enumerate]{itemsep=0.45em, topsep=0.35em, leftmargin=2.4em}
\setlist[itemize]{itemsep=0.25em, topsep=0.25em, leftmargin=1.9em}

\definecolor{mainblue}{RGB}{36,83,140}
\definecolor{lightblue}{RGB}{241,246,252}
\definecolor{lightgray}{RGB}{246,246,246}
\definecolor{rulegray}{RGB}{110,118,128}

\titleformat{\section}
  {\Large\bfseries\color{mainblue}}
  {\thesection.}
  {0.45em}
  {}
  [\vspace{0.2em}\color{rulegray}\titlerule]
\titlespacing*{\section}{0pt}{1.1em}{0.7em}

\pagestyle{fancy}
\fancyhf{}
\fancyhead[L]{\small 错位排列周测 B 卷}
\fancyhead[C]{\small 组合专题：递推、容斥与综合变式}
\fancyhead[R]{\small 刘文博}
\fancyfoot[C]{\small 第 \thepage 页 / 共 \pageref{LastPage} 页}
\renewcommand{\headrulewidth}{0.55pt}
\renewcommand{\footrulewidth}{0.35pt}

\newtcolorbox{infobox}[1]{
  breakable,
  colback=lightblue,
  colframe=mainblue,
  boxrule=0.7pt,
  arc=1mm,
  title=\textbf{#1},
  fonttitle=\bfseries
}

\newenvironment{solution}{\par\noindent\textbf{参考解答：}}{\hfill$\square$\par}
\newcommand{\answerspace}{\vspace{2.3cm}}
\newcommand{\biganswerspace}{\vspace{3.2cm}}
\newcommand{\Dn}{D_n}

\begin{document}

\thispagestyle{empty}
\begin{center}
{\fontsize{22}{28}\selectfont\bfseries 错位排列周测 B 卷\par}
\vspace{0.4cm}
{\large 适合小学高年级至初中低年级数学竞赛训练\par}
\end{center}

\begin{infobox}{考试说明}
\begin{itemize}
  \item 测试时间：75 分钟
  \item 满分：100 分
  \item B 卷更强调方法迁移，请特别注意区分：递推构造、补集法、容斥原理和 $\binom{n}{k}D_{n-k}$ 转化。
\end{itemize}
\end{infobox}

\section{试题}

\begin{enumerate}
  \item （12 分）证明错位排列满足递推公式
  \[
  D_n=(n-1)(D_{n-1}+D_{n-2}).
  \]
  \biganswerspace

  \item （12 分）用容斥原理计算 $D_6$。
  \answerspace

  \item （12 分）8 个不同数字排成一行，恰好有 1 个数字在原位，共有多少种排法？
  \answerspace

  \item （12 分）7 封不同的信放入 7 个对应信封中，至少有 2 封装对，共有多少种装法？
  \answerspace

  \item （12 分）6 个不同数字排成一行，至多有 1 个数字在原位，共有多少种排法？
  \answerspace

  \item （12 分）7 个对象排成一行，要求前 5 个对象都不在原位，后 2 个对象没有限制。共有多少种排法？
  \biganswerspace

  \item （14 分）8 个人排座位，恰好有 2 个人没有坐回自己的座位，共有多少种排法？
  \answerspace

  \item （14 分）8 个不同数字排成一行，至少有 3 个数字在原位，共有多少种排法？
  \biganswerspace
\end{enumerate}

\newpage
\section{详细解答}

\begin{enumerate}
  \item \begin{solution}
  我们只看对象 1 最后去了哪里。

  因为对象 1 不能回到第 1 个位置，所以它只能去其余 $n-1$ 个位置中的某一个，先有
  \[
  n-1
  \]
  种选择。

  设对象 1 去了第 $k$ 个位置（$k\neq 1$），再看原来在第 $k$ 个位置上的对象 $k$：

  \textbf{第一种情况：对象 $k$ 去了第 1 个位置。}  
  这时对象 1 和对象 $k$ 形成一对交换，剩下的 $n-2$ 个对象仍需全部错位排列，因此有
  \[
  D_{n-2}
  \]
  种。

  \textbf{第二种情况：对象 $k$ 没去第 1 个位置。}  
  这时把“第 1 个位置”和“对象 $k$ 原来的位置”一起看作已经处理过，可以把剩下问题转化为 $n-1$ 个对象的错位排列，因此有
  \[
  D_{n-1}
  \]
  种。

  对于对象 1 的每一种去向，上面两类情况总共有
  \[
  D_{n-2}+D_{n-1}
  \]
  种。

  再乘上对象 1 的 $n-1$ 种去向，得到
  \[
  D_n=(n-1)(D_{n-1}+D_{n-2}).
  \]

  证毕。
  \end{solution}

  \item \begin{solution}
  设 $A_i$ 表示“第 $i$ 个对象在原位”，其中 $i=1,2,3,4,5,6$。

  我们要求的是“没有任何一个对象在原位”的排列数，即
  \[
  D_6=6!-\sum |A_i|+\sum |A_i\cap A_j|-\sum |A_i\cap A_j\cap A_k|+\cdots.
  \]

  逐项来看：
  \begin{itemize}
    \item 总排法：$6!=720$；
    \item 固定 1 个对象在原位时，有 $5!$ 种，共 $\binom61 5!=6\times 120=720$；
    \item 固定 2 个对象在原位时，有 $4!$ 种，共 $\binom62 4!=15\times 24=360$；
    \item 固定 3 个对象在原位时，有 $3!$ 种，共 $\binom63 3!=20\times 6=120$；
    \item 固定 4 个对象在原位时，有 $2!$ 种，共 $\binom64 2!=15\times 2=30$；
    \item 固定 5 个对象在原位时，有 $1!$ 种，共 $\binom65 1!=6$；
    \item 固定 6 个对象在原位时，有 $0!=1$ 种，共 $\binom66 0!=1$。
  \end{itemize}

  所以
  \[
  D_6=720-720+360-120+30-6+1.
  \]

  计算得
  \[
  D_6=265.
  \]
  \end{solution}

  \item \begin{solution}
  “恰好有 1 个数字在原位”时，先选出这个在原位的数字，再让剩下 7 个数字全部错位排列。

  因此总数为
  \[
  \binom81 D_7.
  \]

  由递推或已知可得
  \[
  D_7=1854.
  \]

  所以答案是
  \[
  \binom81 D_7=8\times 1854=14832.
  \]
  \end{solution}

  \item \begin{solution}
  “至少有 2 封装对”可以按“恰好有几封装对”分类。

  \textbf{恰好有 2 封装对：}
  \[
  \binom72 D_5=21\times 44=924.
  \]

  \textbf{恰好有 3 封装对：}
  \[
  \binom73 D_4=35\times 9=315.
  \]

  \textbf{恰好有 4 封装对：}
  \[
  \binom74 D_3=35\times 2=70.
  \]

  \textbf{恰好有 5 封装对：}
  先选出哪 5 封装对，有
  \[
  \binom75 D_2=21\times 1=21.
  \]

  \textbf{恰好有 6 封装对：}
  不可能。因为只要已有 6 封装对，最后 1 封也会被迫装对。

  \textbf{恰好有 7 封装对：}
  有
  \[
  1
  \]
  种。

  所以至少有 2 封装对的总数为
  \[
  924+315+70+21+1=1331.
  \]
  \end{solution}

  \item \begin{solution}
  “至多有 1 个数字在原位”包括两类：
  \begin{itemize}
    \item 没有一个在原位；
    \item 恰好有 1 个在原位。
  \end{itemize}

  第一类有
  \[
  D_6=265
  \]
  种。

  第二类：先选哪 1 个在原位，有
  \[
  \binom61=6
  \]
  种；剩下 5 个全部错位排列，有
  \[
  D_5=44
  \]
  种。

  所以第二类有
  \[
  6\times 44=264
  \]
  种。

  总数为
  \[
  265+264=529.
  \]
  \end{solution}

  \item \begin{solution}
  这是“只有部分对象不能回原位”的题，应对前 5 个对象的坏事件使用容斥。

  设 $A_i$ 表示“前 5 个对象中的第 $i$ 个回到原位”，其中 $i=1,2,3,4,5$。那么满足要求的排列数为
  \[
  7!-\binom51 6!+\binom52 5!-\binom53 4!+\binom54 3!-\binom55 2!.
  \]

  逐项计算：
  \[
  7!=5040,
  \]
  \[
  \binom51 6!=5\times 720=3600,
  \]
  \[
  \binom52 5!=10\times 120=1200,
  \]
  \[
  \binom53 4!=10\times 24=240,
  \]
  \[
  \binom54 3!=5\times 6=30,
  \]
  \[
  \binom55 2!=2.
  \]

  所以所求为
  \[
  5040-3600+1200-240+30-2=2428.
  \]
  \end{solution}

  \item \begin{solution}
  如果恰好有 2 个人没有坐回自己的座位，那么另外 6 个人都在原位。

  这 2 个没有回原位的人既不能坐自己的座位，又只能坐剩下的两个空位，所以他们唯一的坐法就是互换座位。

  因此只需先选出是哪 2 个人没有坐回原位，有
  \[
  \binom82=28
  \]
  种；选定以后，他们的错位方式只有 1 种。

  所以总数为
  \[
  28.
  \]
  \end{solution}

  \item \begin{solution}
  “至少有 3 个数字在原位”可以按“恰好有几个数字在原位”分类。

  \textbf{恰好有 3 个在原位：}
  \[
  \binom83 D_5=56\times 44=2464.
  \]

  \textbf{恰好有 4 个在原位：}
  \[
  \binom84 D_4=70\times 9=630.
  \]

  \textbf{恰好有 5 个在原位：}
  \[
  \binom85 D_3=56\times 2=112.
  \]

  \textbf{恰好有 6 个在原位：}
  \[
  \binom86 D_2=28\times 1=28.
  \]

  \textbf{恰好有 7 个在原位：}
  不可能。

  \textbf{恰好有 8 个在原位：}
  有
  \[
  1
  \]
  种。

  因此总数为
  \[
  2464+630+112+28+1=3235.
  \]
  \end{solution}
\end{enumerate}

\end{document}
